\chapter{Состав и реализация модуля, который семантически сопоставляет вакансии и резюме, рассчитывает оценку соответствия в доменном векторном пространстве и ранжирует кандидатов на вакансию}
\label{chapter:impl}
\section{Программный контур подготовки данных и конвейер машинного обучения}
Программная реализация конвейера машинного обучения разделена на офлайн-часть и онлайн-часть в соответствии с описанной для индустриальных систем поиска талантов трехэтапной схемой: подготовка данных о кандидатах, предварительное ранжирование списка соискателей и окончательное переранжирование \cite{src17}. Офлайн-модули собирают, очищают и размечают архивные данные, формируют обучающие пары текст вакансии и текст резюме, и проводят обучение модели, то есть настройку параметров модели на примерах, где значение 1 означает принятие кандидата к рассмотрению или приглашение на следующий этап, а значение 0 означает отказ по конкретной вакансии. Онлайн-часть выполняет вывод модели: преобразует новые тексты вакансий и резюме в плотные векторные представления, рассчитывает оценку их семантического соответствия в доменном векторном пространстве и предоставляет HTTP-маршруты программного интерфейса для ранжирования. Совокупный программный комплекс модуля семантического ранжирования состоит из 18 взаимосвязанных скриптов на языке Python, исполняемых в строго определенном порядке нумерации: от модулей нормализации и извлечения сущностей (скрипты 05-07), через модули подготовки данных для обучения и обучения модели раздельного кодирования (bi-encoder) (08-09), ядро ранжирования цифровых профилей соискателей и сервис программного интерфейса (10-11), который нужен для вызова модели из внешней системы, до серии экспериментальных контрольных запусков конвейера (12-17) и отдельного контура обучения перекрестной модели совместного кодирования (cross-encoder) (18). Все зависимости фиксированы через библиотеки языка Python (sentence-transformers, transformers, scikit-learn, fastapi, uvicorn), что обеспечивает полную воспроизводимость результатов.\par
Под этим конвейером далее понимается набор конкретных подпрограмм, которые последовательно выполняют следующие действия:\par
\begin{compactenum}
  \item Систематический сбор текстов вакансия и резюме и их деидентификация.
  \item Автоматическое отслеживание профильных статистик, например процента дубликатов текстов и плотности профессиональных терминов.
  \item Формирование разметки корпуса данных из системы автоматизированного подбора персонала в рамках подготовки к обучению алгоритма. Данный этап опирается на экспертную разметку в сфере подбора персонала и нужен для формирования набора данных, отражающего особенности этой предметной области \cite{src24}. В будущем планируется адаптировать эти размеченные данные для обучения алгоритмов ранжирования (Learning to Rank).
\end{compactenum}
\section{Реализация и доменная тонкая настройка модели семантического сопоставления текст вакансии и текст резюме}
В реализации основного масштабируемого подхода используется архитектура раздельного кодирования текстов вакансии и резюме (bi-encoder) на базе модели \path|cointegrated/rubert-tiny2|: тексты вакансии и резюме преобразуются в отдельные плотные векторные представления, а затем сравниваются по косинусной близости этих представлений. Это реализовано обучающим скриптом на языке Python \path|experiments/scripts/09_train_biencoder.py|, где консольный аргумент скрипта \path|--model_name| имеет значение по умолчанию \path|cointegrated/rubert-tiny2|, а также сохраненным артефактом \path|experiments/models/bi_encoder_rubert_tiny2|. Под артефактом здесь понимается директория с весами модели, токенизатором, конфигурационными файлами и описанием конвейера обработки, которую можно загрузить для повторного вывода модели без нового обучения; в README.md артефакта указан base\_model: \path|cointegrated/rubert-tiny2|. Для сравнительного эксперимента дополнительно предусмотрена перекрестная модель совместного кодирования (cross-encoder), прошедшая тонкую настройку на той же базовой модели. Этот контур опирается на библиотеку Transformers: в работе Wolf и соавторов она описана как открытая библиотека с единым программным интерфейсом для архитектур Transformer и предобученных моделей, а в коде данной ВКР эта идея используется через класс AutoModelForSequenceClassification для классификации входной последовательности или пары последовательностей \cite{src47}. Скрипт \path|experiments/scripts/18_train_cross_encoder.py| обучает этот класс предсказывать класс кадрового решения для текстовой пары [vacancy\_text, resume\_text] на данных файла train.tsv, валидирует на данных файла val.tsv и сохраняет финальные тестовые метрики только после оценки на данных файла test.tsv.\par
Сохраненный после повторного запуска артефакт оформлен как модель Sentence Transformer: сначала модуль Transformer на базе BertModel строит контекстные представления токенов, затем Pooling агрегирует их в один вектор текста, а Normalize нормирует итоговое векторное представление для устойчивого расчета косинусного сходства. Эта структура соответствует документации Sentence Transformers, где модель описывается как последовательность модулей и отдельно выделяется типовая комбинация Transformer и Pooling с возможным добавлением модуля Normalize; файл \path|modules.json| в артефакте фиксирует именно такую последовательность загрузки модулей \cite{src50}. Векторное сравнение задается через косинусное сходство (cosine similarity). Параметры, зафиксированные в файлах артефакта после повторного обучения 10.04.2026, приведены в таблице.\par
\Needspace{10\baselineskip}
\begingroup
\setlength{\tabcolsep}{3pt}
\renewcommand{\arraystretch}{1}
{\linespread{1}\selectfont
\begin{longtable}{|P{0.30\textwidth}|P{0.34\textwidth}|P{0.28\textwidth}|}
\caption{Параметры сохранённой модели cointegrated/rubert-tiny2}\label{tbl:model-params}\\*
\hline
Параметр & Значение в реализации & Источник \\
\hline
\endfirsthead
\caption[]{Параметры сохранённой модели cointegrated/rubert-tiny2 (продолжение)}\\*
\hline
Параметр & Значение в реализации & Источник \\
\hline
\endhead
\hline
\endfoot
\hline
\endlastfoot
Базовая модель & \path|cointegrated/rubert-tiny2| & \path|09_train_biencoder.py|, README.md артефакта \\
\hline
Тип сохраненной модели & Sentence Transformer & \path|config_sentence_transformers.json|, README.md артефакта \\
\hline
Архитектура модели с механизмом самовнимания & BertModel, model\_type: bert & config.json \\
\hline
Число скрытых слоев & 3 & config.json \\
\hline
Размер скрытого состояния / векторного представления & 312 & config.json, \path|1_Pooling/config.json|, README.md артефакта \\
\hline
Число голов механизма внимания & 12 & config.json \\
\hline
Размер промежуточного слоя & 600 & config.json \\
\hline
Размер словаря & 83828 & config.json \\
\hline
Максимальная длина последовательности & 2048 токенов & README.md артефакта, config.json (max\_position\_embeddings) \\
\hline
Pooling & cls & \path|1_Pooling/config.json| \\
\hline
Нормализация векторного представления & Normalize() & \path|modules.json|, README.md артефакта \\
\hline
Функция сходства & cosine & \path|config_sentence_transformers.json|, README.md артефакта \\
\hline
Файл весов & \path|model.safetensors|, 116,8 МБ & файловая система артефакта после запуска \\
\end{longtable}
}
\endgroup
Обучение выполняется как контрастивная доменная настройка модели раздельного кодирования на текстовых парах текст резюме и текст вакансии: модель получает пары с положительным кадровым решением и пары с отказом из событий подбора персонала и учится сближать векторные представления пар, по которым специалист по подбору персонала перевел кандидата к дальнейшему рассмотрению, одновременно отдаляя пары, завершившиеся отказом. Под пространством векторных представлений здесь понимается числовое пространство, в котором каждый текст представлен вектором, а близость векторов отражает близость смыслов. После подготовки данных сэмплирование выполняется на основе связанных с резюме и вакансиями событий подбора персонала: скрипт \path|experiments/scripts/01_sample_and_profile_data.py| читает первоисточник данных по пути \path|data/_/dwh_dwh_hr_received_resume_event_ods.tsv|, отбирает только явно интерпретируемые статусы, например отказы, приглашения, статусы рассмотрения, собеседования или принятия, и затем потоково извлекает соответствующие строки из сырых таблиц базы данных, характеризующих сущности резюме и вакансий.\par
Разметка строится консервативно: класс отказа фиксируется для явных отказов (Отклонено, Отклоненные, текстовые формы отказ*), а положительный класс присваивается для явных статусов рассмотрения, приглашения, собеседования или принятия. Неоднозначные события подбора персонала, например промежуточные статусы без финального решения, технические изменения карточки, повторные просмотры или статусы, по которым невозможно понять, было ли принято кадровое решение, не используются для обучения. Роль специалистов по подбору персонала в реализованном контуре ограничена происхождением меток в журнале событий; отдельного участия специалистов в выводе модели, программном интерфейсе или ручной корректировке результатов ранжирования в коде не предусмотрено.\par
Фактическая пересборка данных дала следующие числа: из 1 756 716 просмотренных событий из нашего набора данных найдено 114 145 явно размеченных событий (принять или отказать) и 112 698 уникальных явно размеченных сущностей резюме-вакансия. Для обучающего сэмпла выбрано 10 000 пар, что обусловлено вычислительными и временными ограничениями; после проверки наличия текстов резюме и вакансий осталось 5 218 связанных событий. Производные файлы содержат 7 098 резюме, 3 262 вакансии и 5 218 текстовых текст резюме и текст вакансии. Файлы обучения, валидации и тестирования после стратифицированного разбиения содержат соответственно файл train.tsv - 4 176 пар, файл val.tsv - 521 пару, файл test.tsv - 521 пару. Распределение меток: train - 1 423 пары, где кандидат был принят к рассмотрению или приглашен на следующий этап, и 2 753 пары, где был зафиксирован отказ; val - 177 и 344; test - 177 и 344 соответственно. Проверка связности показала, что все идентификаторы resume\_id и vacancy\_id из файлов присутствуют в файлах таблиц, выгруженных из базы данных \path|resumes_unified.tsv| и \path|vacancies_unified.tsv|. Признаков искажения кодировки, то есть появления нечитаемых последовательностей символов из-за неверной интерпретации кодировки текста, в model\_input не найдено.\par
Непосредственная загрузка обучающих размеченных примеров вакансия-резюме реализована в скрипте \path|experiments/scripts/08_model_dataset.py|. В этом файле класс StreamingPairDataset наследуется от IterableDataset библиотеки PyTorch. Такое решение опирается на модель разработки PyTorch: в работе Paszke и соавторов эта библиотека описана как среда высокопроизводительного глубокого обучения, на которой построен используемый механизм потоковой выдачи обучающих примеров \cite{src46}. Класс потоково читает данные в текстовом формате для хранения таблиц (Tab-Separated Values), получает тексты model\_input из таблицы под названием features базы данных в формате SQLite и на каждой итерации возвращает объект InputExample(texts=[resume\_text, vacancy\_text], label=label) для библиотеки Sentence Transformers. Документация Sentence Transformers связывает обучение модели с набором данных, функцией потерь и числовой меткой пары; в реализованном варианте такая метка определяет, должна ли функция потерь сближать или разводить векторные представления резюме и вакансии \cite{src50}. После исправления текста model\_input для резюме строится из неперсональных полей: желаемая должность desired\_profession, достижения best, сведения о себе dop, уровень владения компьютером computer. ФИО, телефоны, email и адреса не используются, поскольку являются персональными данными. Для вакансий текстовая часть строится из фактических полей описания требований к кандидату candidat и описания компании company, пустое устаревшее поле info не используется.\par
Гиперпараметры и фактические числа повторного запуска обучения приведены в таблице.\par
\Needspace{10\baselineskip}
\begingroup
\setlength{\tabcolsep}{3pt}
\renewcommand{\arraystretch}{1}
{\linespread{1}\selectfont
\begin{longtable}{|P{0.30\textwidth}|P{0.34\textwidth}|P{0.28\textwidth}|}
\caption{Гиперпараметры и фактические показатели обучения}\label{tbl:hyperparams}\\*
\hline
Параметр запуска & Значение & Источник \\
\hline
\endfirsthead
\caption[]{Гиперпараметры и фактические показатели обучения (продолжение)}\\*
\hline
Параметр запуска & Значение & Источник \\
\hline
\endhead
\hline
\endfoot
\hline
\endlastfoot
Python & 3.11.9 & локальное .venv, README.md артефакта \\
\hline
Sentence Transformers & 5.4.0 & .venv, README.md артефакта \\
\hline
Transformers & 5.5.3 & .venv, README.md артефакта \\
\hline
PyTorch & 2.11.0+cpu & .venv, README.md артефакта \\
\hline
Base model & \path|cointegrated/rubert-tiny2| & training command, \path|09_train_biencoder.py| \\
\hline
Dataset loader & StreamingPairDataset + DataLoader & \path|09_train_biencoder.py|, \path|08_model_dataset.py| \\
\hline
Train / val / test pairs & \texttt{4176 / 521 / 521} & фактический подсчет строк в \path|data/splits|/*.tsv \\
\hline
Batch size & 8 & training command, \path|09_train_biencoder.py|, stdout запуска \\
\hline
Epochs & 3 & training command, \path|09_train_biencoder.py|, stdout запуска \\
\hline
Train steps & 1566 & stdout прогресса обучения \\
\hline
Loss & ContrastiveLoss & \path|09_train_biencoder.py|, README.md артефакта \\
\hline
Distance metric для loss & \path|SiameseDistanceMetric.COSINE_DISTANCE| & README.md артефакта \\
\hline
Margin & 0.5 & README.md артефакта \\
\hline
Learning rate & 5e-05 & README.md артефакта \\
\hline
Optimizer & adamw\_torch\_fused & README.md артефакта \\
\hline
Gradient accumulation & 1 & README.md артефакта \\
\hline
\path|FP16/BF16| & False / False & README.md артефакта \\
\hline
Warmup steps & command: 100; generated model card: 0 & training command и README.md артефакта; значение применения требует отдельной проверки Trainer API \\
\hline
Train runtime & 6217 с & stdout успешного запуска \\
\hline
Train samples per second & 2.015 & stdout успешного запуска \\
\hline
Train steps per second & 0.252 & stdout успешного запуска \\
\hline
Train loss & 0.008956 & stdout успешного запуска \\
\hline
Общая длительность команды & 6254.096 с & внешний замер Stopwatch вокруг training command \\
\hline
Validation accuracy & 0.9731 при threshold 0.7408 & финальная оценка Binary Classification Evaluator на val.tsv \\
\hline
Validation F1 & 0.9605 при threshold 0.7093 & финальная оценка Binary Classification Evaluator на val.tsv \\
\hline
Validation precision / полнота & \texttt{0.9605 / 0.9605} & финальная оценка Binary Classification Evaluator на val.tsv \\
\hline
Validation AP / MCC & \texttt{0.9870 / 0.9401} & финальная оценка Binary Classification Evaluator на val.tsv \\
\hline
Output path & \path|experiments/models/bi_encoder_rubert_tiny2| & \path|09_train_biencoder.py|, stdout успешного запуска \\
\end{longtable}
}
\endgroup
Алгоритм ранжирования соискателей реализован в классе ATSRanker файла \path|experiments/scripts/10_ranking_module.py|. Текущая реализация использует последовательный конвейер из трех операций: жесткая фильтрация набора кандидатов, преобразование текста вакансии и резюме в векторные представления через объект Sentence Transformer, то есть загруженную модель кодирования предложений из библиотеки sentence-transformers, затем вычисление косинусного сходства полученных векторных представлений и сортировка результатов фильтрации по убыванию числовой оценки семантического соответствия резюме требованиям вакансии. В будущем планируется реализация методов обучения ранжированию (Learning to Rank), метода попарного обучения ранжированию под названием RankNet, интеграции с векторной базой данных Milvus, алгоритмами приближенного поиска ближайших соседей FAISS, HNSW, DiskANN, SPANN, отдельного индекса приближенного поиска ближайших соседей (Approximate Nearest Neighbor, ANN), а также модуля справедливого ранжирования, который будет реализовывать логику оценки и снижения несправедливости по отношению к кандидатам со стороны алгоритмов ранжирования \cite{src12, src13, src14, src15, src32, src33, src36}.\par
Жесткая фильтрация выполняется методом \_apply\_hard\_filters(candidates, filters). Если фильтры не переданы, набор кандидатов возвращается без изменений. Если фильтры заданы, для каждого кандидата проверяется точное совпадение значения в карточке кандидата и требуемого значения из запроса: выражение cand.get(k) == required\_val означает, что по ключу k, например city, grade или specialization, значение в словаре кандидата должно совпасть со значением фильтра required\_val. Под словарем filters понимается структура вида \texttt{\{"city": "Москва", "grade": "Middle"\}}, переданная в запросе к программному интерфейсу. Записи кандидатов, не удовлетворяющие хотя бы одному условию, заранее удаляются из списка, который допускается до дальнейшего рассмотрения. После этого нейросетевая модель сравнивает с вакансией только оставшихся кандидатов, то есть рассчитывает не абстрактный единый балл, а доменную числовую оценку соответствия между требованиями вакансии и подготовленным текстом резюме с учетом тонкой настройки модели на кадровых событиях. Это не поиск по индексу и не обращение к внешней базе данных, а прямой проход по списку объектов-кандидатов, хранящихся в оперативной памяти процесса Python.\par
Семантический этап реализован в методе search. Этот метод получает текст вакансии vacancy\_text, набор кандидатов candidate\_pool, необязательные фильтры filters и ограничение top\_k. После фильтрации модуль извлекает поле model\_input из оставшихся кандидатов, преобразует текст вакансии и тексты кандидатов в векторные представления на лету, вычисляет \path|util.cos_sim|, выбирает top\_k через функцию topk библиотеки torch и формирует ответ в формате объектной нотации JavaScript. Близкий к коду фрагмент реализации:\par
{\linespread{1}\selectfont
\begin{lstlisting}[language=Python, numbers=left, stepnumber=1, caption={Семантический этап ранжирования в ATSRanker}, label={lst:search}]
filtered_candidates = self._apply_hard_filters(candidate_pool, filters)
candidate_texts = [c.get('model_input', '') for c in filtered_candidates]
vacancy_emb = self.model.encode(vacancy_text, convert_to_tensor=True)
candidate_embs = self.model.encode(candidate_texts, convert_to_tensor=True)
cos_scores = util.cos_sim(vacancy_emb, candidate_embs)[0]
top_results = torch.topk(cos_scores, k=min(top_k, len(filtered_candidates)))
\end{lstlisting}
}
Возвращаемый результат имеет машиночитаемый формат объектной нотации JavaScript: поле status показывает успешность выполнения операции, latency\_ms содержит измеренную длительность обработки запроса в миллисекундах, filters\_applied фиксирует примененные жесткие фильтры, candidates\_evaluated показывает число кандидатов после фильтрации, а массив results содержит итоговую выдачу. Для каждого элемента выдачи сохраняются candidate\_id, то есть идентификатор кандидата, округленный match\_score, то есть числовая оценка семантического соответствия резюме требованиям вакансии, и короткая текстовая summary, сформированная из начала поля model\_input. Поле latency\_ms вычисляется внутри метода search, однако контролируемый запуск по оценке временных показателей работы модуля ранжирования на момент выполнения работы не запускался, поэтому численные утверждения о задержке, пропускной способности, памяти или масштабировании в текущем разделе не приводятся.\par
Интеграция модулем ранжирования через программный интерфейс реализована в файле \path|experiments/scripts/11_api_server.py|. Сервер на фреймворке языка Python под названием FastAPI хранит кандидатские профили во временном списке в оперативной памяти (in memory) CANDIDATE\_INDEX: маршрут POST \path|/index/candidate| добавляет объект с идентификатором id, подготовленным текстом model\_input и необязательными метаданными; обработчик маршрута POST \path|/search/rank| вызывает \path|RANKER_ENGINE.search|(vacancy\_text, CANDIDATE\_INDEX, filters, top\_k) и формирует JSONResponse из результата; маршрут POST \path|/score| считает одиночный match\_score для пары текст вакансии и текст резюме. Следовательно, реализованная интеграция является прототипом программного интерфейса с хранением данных в оперативной памяти (in memory), а не промышленной связкой с Milvus, DiskANN, SPANN, FAISS, HNSW или другой векторной базой данных \cite{src12, src13, src14, src32, src33}.\par
Оценка качества текущей модели раздельного кодирования (bi-encoder) интерпретируется как проверка семантического расчета соответствия резюме требованиям вакансии с учетом тонкой настройки на сформированных файлах выборок, а не как внешняя бизнес-валидация системы автоматизированного подбора персонала. Для файлов val.tsv и test.tsv были измерены метрики бинарного разделения по косинусной оценке близости с подбором порога на валидационном срезе данных. После тонкой настройки (fine-tuning) сохраненная модель \path|experiments/models/bi_encoder_rubert_tiny2| достигла на валидационной подвыборке средней точности (Average Precision) 0.9870, лучшей гармонической средней между точностью и полнотой (F1-score) 0.9582, лучшей доли правильных ответов (Accuracy) 0.9712 и коэффициента корреляции Мэтьюса (Matthews Correlation Coefficient, MCC) 0.9365. На test выборке модель показала среднюю точность 0.9921, лучшую гармоническую среднюю между точностью и полнотой 0.9636, лучшую долю правильных ответов 0.9750 и MCC 0.9447. Эти значения показывают, что после дообучения модель раздельного кодирования формирует распределение оценок, при котором пары, где специалист по подбору персонала продвинул кандидата к дальнейшему рассмотрению, отделяются от пар, где был зафиксирован отказ, на уровне средней точности 0.9921, гармонической средней (F1-score) 0.9636, доли правильных ответов 0.9750 и коэффициента корреляции Мэтьюса 0.9447 на тестовой выборке. Все эти метрики не являются метриками справедливого ранжирования результатов (fair ranking) или внешней экспертной проверкой; они показывают качество ранжирования на наших размеченных данных при условии, что разметка корректно отражает решения специалистов по подбору персонала.\par
Для проверки метрик ранжирования был выполнен отдельный полный прогон по файлу тестовой выборки \path|data/splits/test.tsv| с группировкой по идентификатору вакансии vacancy\_id: внутри каждой вакансии резюме ранжировались по косинусной оценке близости текущей модели после тонкой настройки (fine-tuning). Необработанный вывод сохранен в файл \path|experiments/results/full_test_ranking_metrics_stdout.txt|, агрегированные результаты и значения на уровне групп вакансий - в файл \path|experiments/results/full_test_ranking_metrics.json|. В файле test.tsv содержится 521 пара, 423 уникальные вакансии, 510 уникальных резюме, 177 пар, где кандидат был принят к рассмотрению или приглашен на следующий этап, и 344 пары, где по отклику был зафиксирован отказ. Пропущенных идентификаторов resume\_id и vacancy\_id относительно объединенных файлов нет. При этом тестовая выборка разрежена: для 304 вакансий нет ни одного кандидата, которого специалист по подбору персонала исторически продвинул на следующий этап, только 59 вакансий имеют минимум двух кандидатов, оставивших отклик, а нормализованная дисконтированная кумулятивная выгода (Normalized Discounted Cumulative Gain, NDCG) математически определена только для 33 вакансий, где есть хотя бы два кандидата и хотя бы один кандидат с зафиксированным продвижением.\par
\Needspace{10\baselineskip}
\begingroup
\setlength{\tabcolsep}{3pt}
\renewcommand{\arraystretch}{1}
{\linespread{1}\selectfont
\begin{longtable}{|P{0.30\textwidth}|P{0.34\textwidth}|P{0.28\textwidth}|}
\caption{Ranking-метрики на полном test split}\label{tbl:ranking}\\*
\hline
Метрика на полном "test.tsv" & Значение & Область усреднения \\
\hline
\endfirsthead
\caption[]{Ranking-метрики на полном test split (продолжение)}\\*
\hline
Метрика на полном "test.tsv" & Значение & Область усреднения \\
\hline
\endhead
\hline
\endfoot
\hline
\endlastfoot
NDCG@1 & 1.0000 & 33 вакансии, пригодные для NDCG \\
\hline
NDCG@3 & 1.0000 & 33 вакансии, пригодные для NDCG \\
\hline
NDCG@5 & 1.0000 & 33 вакансии, пригодные для NDCG \\
\hline
Precision@1 & 1.0000 & 119 вакансий с хотя бы одним исторически продвинутым кандидатом \\
\hline
Precision@3 & 0.9958 & 119 вакансий с хотя бы одним исторически продвинутым кандидатом \\
\hline
Precision@5 & 0.9958 & 119 вакансий с хотя бы одним исторически продвинутым кандидатом \\
\hline
Recall@1 & 0.8471 & 119 вакансий с хотя бы одним исторически продвинутым кандидатом \\
\hline
Recall@3 & 0.9824 & 119 вакансий с хотя бы одним исторически продвинутым кандидатом \\
\hline
Recall@5 & 0.9924 & 119 вакансий с хотя бы одним исторически продвинутым кандидатом \\
\hline
MRR & 1.0000 & 119 вакансий с хотя бы одним исторически продвинутым кандидатом \\
\end{longtable}
}
\endgroup
Если усредлять точность на первых K позициях (Precision at K, Precision@K) по всем 423 вакансиям, включая 304 вакансии, для которых в тестовой разметке нет ни одного кандидата, продвинутого специалистом по подбору персонала на следующий этап, то Precision@1 составляет 0.2813, Precision@3 - 0.2801, Precision@5 - 0.2801. Полнота (Recall) и нормализованная дисконтированная кумулятивная выгода (Normalized Discounted Cumulative Gain, NDCG) для таких групп не определены по смыслу, потому что в них отсутствует кандидат, которого историческое решение специалиста перевело на следующий этап и которого модель могла бы поднять в выдаче. Поэтому в тексте работы основные метрики ранжирования приводятся для вакансий с хотя бы одним исторически продвинутым кандидатом, а разреженность тестового набора фиксируется как ограничение эксперимента.\par
Как проектное допущение дальнейшего развития, элементы, не реализованные в текущем коде, фиксируются как направления будущей работы: приближенный поиск ближайших соседей и внешнее векторное хранилище для больших наборов кандидатов; переранжирование на основе обучения ранжированию (Learning to Rank), включая метод попарного обучения ранжированию под названием RankNet или списочные функции потерь; более строгая оценка по нормализованной дисконтированной кумулятивной выгоде (Normalized Discounted Cumulative Gain, NDCG), точности (Precision), полноте (Recall) и средней обратной позиции первого исторически продвинутого кандидата (Mean Reciprocal Rank, MRR) на специально сформированных списках выдачи с несколькими кандидатами на каждую вакансию; аудит предвзятости и механизмы справедливого ранжирования результатов (fair ranking); контролируемое нагрузочное тестирование с фиксированным окружением, размером набора кандидатов и аппаратной конфигурацией.\par
\section{Выводы по третьему разделу}
\begin{compactenum}
  \item Спроектированный и реализованный программный контур включает подготовку обучающих пар, обучение модели раздельного кодирования, сохранение артефакта модели, класс ATSRanker и сервисный API на базе FastAPI.
  \item Доменная тонкая настройка модели \path|cointegrated/rubert-tiny2| позволяет связать семантическое расстояние между текстами вакансии и резюме с историческими кадровыми событиями, а не только с общеязыковой близостью текстов.
  \item Класс ATSRanker инкапсулирует практическую логику ранжирования: подготовку входов, получение векторных представлений, расчет косинусной близости и сортировку кандидатов по итоговой оценке.
  \item Полученные ranking-метрики показывают работоспособность прототипа на отложенной выборке, но разреженность групп вакансий требует аккуратной интерпретации результата и дальнейшей проверки на более плотных списках кандидатов.
\end{compactenum}
